\chapter{Random Graphs, Derandomization, and Symmetry Breaking}
\label{ch:stage5-random-graphs}

This supplement develops the random-structure and derandomization ideas that
connect the probabilistic method to algorithms.  Existing chapters retain
their treatments of the Local Lemma, expanders, conditional expectations, and
parallel MIS; this chapter supplies the missing bridges and student-level
threshold calculations.

\section{The Erd\H{o}s--R\'enyi models}

In $G(n,p)$, every one of the $\binom n2$ possible edges is present
independently with probability $p$.  In $G(n,m)$, a uniformly random set of
$m$ edges is chosen.  Many monotone properties have comparable thresholds in
the two models, but a proof must identify the coupling or concentration step
that transfers a result.

\begin{proposition}[Connectivity threshold]
If $p=(\log n+c)/n$ for fixed $c$, then
\[
 \Pr[G(n,p)\text{ is connected}]\longrightarrow e^{-e^{-c}}.
\]
In particular, the threshold for connectivity is asymptotically
$(\log n)/n$.
\end{proposition}

The main obstruction near this threshold is an isolated vertex.  The expected
number of isolated vertices is
$n(1-p)^{n-1}\sim e^{-c}$, and a second-moment argument controls the
probability that several such obstructions occur.

\section{Hamilton cycles}

A Hamilton cycle is a cycle visiting every vertex exactly once.  The sharp
threshold in $G(n,p)$ is
\[
 p=\frac{\log n+\log\log n+\omega(1)}{n},
\]
where $\omega(1)\to\infty$: above this threshold the graph is Hamiltonian
with high probability, while below the corresponding threshold it usually
has a vertex of degree at most one and cannot be Hamiltonian.

This example teaches a general threshold workflow: identify a local
obstruction for the lower bound, count candidate structures for the upper
bound, and then use a sprinkling or rotation-extension argument to connect
the surviving structure globally.

\section{Second-moment method for random structures}

Let $X$ count copies of a target structure, such as triangles or Hamilton
cycles.  The first moment proves nonexistence when $\mathbb{E}X=o(1)$.  The
second moment can prove existence:
\[
 \Pr[X>0]\geq\frac{(\mathbb{E}X)^2}{\mathbb{E}[X^2]}.
\]
Expanding $\mathbb{E}[X^2]$ by overlap type is often the main technical step.
If $\operatorname{Var}(X)=o((\mathbb{E}X)^2)$, then $X>0$ with probability
$1-o(1)$ and $X/\mathbb{E}X\to1$ in probability.

\section{Ramsey bounds}

Color each edge of $K_n$ red or blue independently.  For a fixed $k$-set,
the probability that it is monochromatic is
$2^{1-\binom{k}{2}}$.  Hence the expected number of monochromatic $K_k$s is
\[
 \binom nk 2^{1-\binom{k}{2}}.
\]
If this quantity is less than one, some coloring has no monochromatic
$K_k$, proving a lower bound on the diagonal Ramsey number.  The argument is
an existence proof; it does not by itself provide an efficient coloring
algorithm.

\section{Expander codes}

Let $G=(L,R,E)$ be a left-regular bipartite graph.  Associate one message
bit to each left vertex and one parity constraint to each right vertex.  A
good expansion property lets local parity checks detect and correct sparse
errors.

\begin{theorem}[Expander-code decoding, qualitative form]
Suppose every set $S\subseteq L$ of size at most $s$ has at least
$(1-\varepsilon)d|S|$ distinct neighbors, where $d$ is the left degree and
$\varepsilon$ is a sufficiently small constant.  Then the associated
expander code has linear minimum distance for an appropriate constraint
construction, and a bit-flipping decoder corrects a constant fraction of
errors in linear time.
\end{theorem}

The constants depend on the precise code and decoder.  The proof counts
unique neighboring checks: if a small error set had too many unsatisfied
checks hidden by cancellations, it would contradict expansion.  Random
regular bipartite graphs supply suitable expanders with high probability, but
the construction and decoding theorem must specify the degree and expansion
parameters.

\section{Constructive Local Lemma}

In the variable model, each bad event depends on a subset of independent
variables.  Moser--Tardos repeatedly resamples the variables of a currently
true bad event.  The witness-tree proof associates each resampling with a
tree whose labels record dependency edges.  Under the asymmetric LLL
condition
\[
 \Pr[A_i]\leq x_i\prod_{j\in\Gamma(i)}(1-x_j),
\]
the expected number of resamplings of $A_i$ is at most $x_i/(1-x_i)$.

The entropic viewpoint reaches the same conclusion by showing that a long
resampling log would encode too much information relative to the entropy of
the fresh random variables.  It is useful pedagogically, but neither proof
implies an efficient implementation unless bad events can be detected and
resampled efficiently.

\section{Bit fixing and conditional expectations}

Suppose random bits $X_1,\ldots,X_n$ define a nonnegative objective $Z$ and
$\mathbb{E}Z\leq B$.  At step $i$, choose the value of $X_i$ that keeps
$\mathbb{E}[Z\mid X_1,\ldots,X_i]\leq B$.  The tower property guarantees that
at least one of the two choices has this property.  After all bits are fixed,
the resulting deterministic outcome satisfies $Z\leq B$.

The method is algorithmic only when conditional expectations can be computed
or bounded efficiently.  This distinction separates a mathematical
derandomization proof from an implementable deterministic algorithm.

\section{Limited-independence derandomization}

Pairwise-independent variables preserve expectations and variances of sums,
but generally do not preserve Chernoff tails.  $k$-wise independence can
preserve moment bounds up to order $k$, and small-bias spaces can replace
uniform bits for some linear tests.  A derandomization argument must match
the pseudorandomness property to the analysis being used.

\section{Symmetry breaking}

When processors begin in identical states, deterministic protocols may be
unable to select different actions.  Random priorities break symmetry.  A
typical round gives each active processor an independent priority, selects
local winners, and removes winners and their conflicts.  The analysis counts
the expected surviving conflicts and proves geometric progress.

Ethernet-style randomized backoff uses the same principle: after a collision,
each participant chooses a random delay.  Correct performance statements
must specify the collision model, feedback assumptions, and whether the
adversary is adaptive.

\section{Student checklist}

For a random-graph or derandomization proof, students should identify:

\begin{enumerate}
\item the random model and its independent choices;
\item the obstruction or witness being counted;
\item whether a first moment, second moment, martingale, or LLL argument is
      appropriate;
\item the exact threshold and asymptotic regime;
\item whether the result is existential, randomized algorithmic, or
      deterministic and efficiently constructible;
\item the implementation assumptions for resampling or symmetry breaking.
\end{enumerate}

\begin{problem}[Connectivity threshold]
Compute the expected number of isolated vertices in $G(n,p)$ and use the
first moment to prove that $p=(\log n-\omega(1))/n$ is below the connectivity
threshold.
\end{problem}

\begin{problem}[Second moment]
For a selected random subgraph count copies of a fixed small graph $H$.
Expand the second moment by overlap size and identify when the variance is
lower order than the squared expectation.
\end{problem}

\begin{problem}[Conditional expectations]
Derandomize a randomized MAX-SAT assignment by explicitly computing the
conditional expected number of satisfied clauses after each variable choice.
\end{problem}

\begin{problem}[Symmetry breaking]
Analyze one round of random-priority MIS and show why the expected number of
remaining edges decreases under a suitable priority rule.
\end{problem}
